\section{Supplementary statistical and systems audit}
\label{app:statistics}

\paragraph{Sources, order, and linkage.}
The author-reported complete dashboard history comprises 1,704 rated entries
from August 24--30: 483 bargaining, 562 negotiation, and 659 persuasion. The
underlying export and parser are unavailable in this package, so these
descriptive summaries cannot be independently regenerated from it. Four games
were against humans and remain in the reported history. The platform starts
each family rating at 1,000. Summed displayed changes are rounded to one
decimal and need not equal the difference between this start and the
higher-precision final rating; platform-wide recalculation can also prevent
exact additivity.

Chronological thirds sum to $+364.4/-27.6/+199.1$ for bargaining,
$+422.8/+159.0/+113.2$ for negotiation, and $+508.1/+261.7/+63.2$ for
persuasion. Long-run role means were $+1.11$ Alice, $+1.11$ Bob, $+0.67$
negotiation buyer, $+1.84$ negotiation seller, $+1.71$ persuasion buyer, and
$+0.86$ persuasion seller. These partitions describe attenuation and role
asymmetry across changing policies and opponents; they are not stationarity
tests or estimates of version effects. The family and role statistics for this
complete history are not machine-readable in the supplied package.

The rating analysis uses all 248 live rows in
\path{data/v33_rating_outcomes.csv} after excluding the explicitly named
synthetic contract-test row. File order is the executed assignment order; the
analysis writes \path{data/v33_rating_outcomes_sequenced.csv} with an explicit
1--248 index. The append-only evidence and session report independently retain
checkpoint order and assigned game IDs.

Table~\ref{tab:leaderboard-snapshot} records the separate author-supplied
August 31 dashboard snapshot; it is not reconstructed from assignment deltas.
\begin{table}[t]
  \caption{Leaderboard snapshot supplied by the author. The player label is
  the study's agent alias.}
  \label{tab:leaderboard-snapshot}
  \centering\scriptsize
  \begin{tabular}{rllrrrr}
    \toprule
    Rank & Player & Type & Rating & Bargaining & Negotiation & Persuasion \\
    \midrule
    95 & ELP-agent & Agent & 1674.3 & 1535.9 & 1661.3 & 1825.8 \\
    \bottomrule
  \end{tabular}
\end{table}

The text export contains 220 complete records: 127/150 bargaining, 74/78
persuasion, and 19/20 negotiation (88.7\% overall). Two records omit a repeated
header but retain separate setup and result blocks, which the parser recovers.
The export has no game IDs, so transcript records are not positionally joined to
ratings. No missing-at-random assumption is made; behavioral quantities are
available-case process summaries.
For example, the 113 observed bargaining agreements imply only the worst-case
full-run range 113/150--136/150 (75.3--90.7\%) without assumptions about the 23
missing records.

\paragraph{Descriptive rating paths.}
Table~\ref{tab:v33-distribution} reports signs, means, medians, dispersion, and
no confidence interval. The positive bargaining mean coexists with an almost
even sign split; persuasion has a positive mean but negative median. Lag-one
sample autocorrelations were $0.05$, $-0.08$, and $0.07$ for bargaining,
persuasion, and negotiation, but do not establish independence. Repeated
opponents cannot be identified from the rating extract. Consequently neither a
fixed-sample IID interval nor a post hoc confidence sequence would have a clear
estimand or justified dependence structure.

Figure~\ref{fig:v33-sequential} plots cumulative changes against the executed
family guards. Bargaining's maximum within-run drawdown was 24.75 points before
the 20-game activation point. Persuasion crossed its 22-point boundary at game
78; negotiation crossed 30 points at game 20.

\begin{figure}[p]
\centering
\definecolor{ratingNavy}{HTML}{24547D}
\definecolor{ratingCoral}{HTML}{C04E56}
\definecolor{ratingPurple}{HTML}{75629A}
\pgfplotsset{rating panel/.style={
  width=.84\textwidth,height=4.45cm,
  axis x line*=bottom,axis y line*=left,
  xlabel={Assignment},ylabel={Rating change},
  tick label style={font=\scriptsize},
  label style={font=\scriptsize},
  title style={font=\footnotesize\bfseries},
  ymajorgrids=true,major grid style={black!12},
  scaled ticks=false,
  clip=true,
  every axis plot/.append style={mark=none}
}}

\begin{tikzpicture}
\begin{axis}[rating panel,title={Bargaining: 150 games, $+190.94$},
  xmin=0,xmax=150,ymin=-40,ymax=210,
  xtick={0,25,50,75,100,125,150},ytick={-30,0,50,100,150,200}]
  \addplot[color=ratingNavy,line width=1pt]
    table[x=game,y=cumulative]{figures/ratings_bargaining.dat};
  \addplot[color=ratingCoral,dashed,line width=.8pt]
    table[x=game,y=trailing]{figures/ratings_bargaining.dat};
  \addplot[color=ratingPurple,densely dotted,line width=.8pt]
    coordinates {(0,-30) (150,-30)};
  \addplot[color=black!45,densely dashed,line width=.6pt]
    coordinates {(20,-40) (20,210)};
\end{axis}
\end{tikzpicture}\par\smallskip

\begin{tikzpicture}
\begin{axis}[rating panel,title={Persuasion: 78 games, $+55.67$},
  xmin=0,xmax=78,ymin=-35,ymax=85,
  xtick={0,20,40,60,78},ytick={-30,-10,10,30,50,70}]
  \addplot[color=ratingNavy,line width=1pt]
    table[x=game,y=cumulative]{figures/ratings_persuasion.dat};
  \addplot[color=ratingCoral,dashed,line width=.8pt]
    table[x=game,y=trailing]{figures/ratings_persuasion.dat};
  \addplot[color=ratingPurple,densely dotted,line width=.8pt]
    coordinates {(0,-28) (78,-28)};
  \addplot[color=black!45,densely dashed,line width=.6pt]
    coordinates {(20,-35) (20,85)};
\end{axis}
\end{tikzpicture}\par\smallskip

\begin{tikzpicture}
\begin{axis}[rating panel,title={Negotiation: 20 games, $-33.73$},
  xmin=0,xmax=20,ymin=-43,ymax=2,
  xtick={0,5,10,15,20},ytick={-40,-30,-20,-10,0}]
  \addplot[color=ratingNavy,line width=1pt]
    table[x=game,y=cumulative]{figures/ratings_negotiation.dat};
  \addplot[color=ratingCoral,dashed,line width=.8pt]
    table[x=game,y=trailing]{figures/ratings_negotiation.dat};
  \addplot[color=ratingPurple,densely dotted,line width=.8pt]
    coordinates {(0,-40) (20,-40)};
  \addplot[color=black!45,densely dashed,line width=.6pt]
    coordinates {(20,-43) (20,2)};
\end{axis}
\end{tikzpicture}

\caption{Cumulative displayed rating changes from the ordered rating CSV.
Solid blue is cumulative change; dashed coral is running peak minus the family
drawdown limit; dotted purple is the baseline-loss limit. The vertical gray
line marks the 20-game guard activation point. These controller boundaries
are not inferential thresholds.}
\label{fig:v33-sequential}
\end{figure}


\paragraph{Support and confounding.}
Table~\ref{tab:v33-persuasion-support} gives the full role--mode--value-visibility
support. Every hidden-value persuasion assignment placed ELP in the seller
role. Hidden-versus-known marginal totals therefore combine information and
role effects and are not compared. The machine-readable context CSV retains
role, information, horizon, discount, signal mode, prior, and cutoff buckets;
small cells remain exploratory and no multiplicity-adjusted claims are made.

\paragraph{Available-case behavior.}
Among 127 exported bargaining games, 113 reached agreement and 14 ended in
walkaway; agreements used a median one observed round (mean 1.40, maximum 12).
Negotiation produced 11/19 agreements, four no-deals, and four walkaways. The 74
persuasion games contain 1,480 decisions. In 40 seller games, 397/397 observed
high-quality products and 135/403 low-quality products received a positive
signal. Opponents bought 309/532 positive-signal and 24/268 negative-signal
products. In 34 buyer games, ELP bought 350/680 products, including 245 high and
105 low; passed quality is hidden. None of these unlinked summaries estimates a
rating effect or counterfactual behavior.

\paragraph{Guard sensitivity.}
Table~\ref{tab:v33-stop-sensitivity} replays the observed sequence while jointly
scaling all family and role thresholds. At half thresholds, bargaining would
have stopped at game 17, persuasion at 24, and negotiation at 19. At $1.25$
times the thresholds, neither adaptively stopped family crosses a guard in its
observed prefix. These are path-dependent replays, not counterfactual estimates:
later matchmaking need not remain the same after a different stopping decision.

\paragraph{Executed horizons and safeguards.}
All 78 persuasion games recorded 20 rounds, so its missing-horizon terminal
branch was never activated. Horizons were unknown in 76 bargaining games (113
decisions) and 10 negotiation games (28 decisions). One complete-information
seller game with an unknown horizon triggered the short-horizon repair on all
four of its decisions; two further repair activations occurred in known
short-horizon games. A future policy should use a separate unknown-horizon
branch rather than encode it as zero.

The evidence-level audit in
\path{data/v33_execution_activation_audit.csv} counts 1,838 live policy
decisions. Table~\ref{tab:v33-activations} distinguishes safeguards that fired
from those merely present in code.
\begin{table}[t]
  \caption{Executed safeguard and repair activations. ``Units'' counts distinct
  games, or families for the drawdown row.}
  \label{tab:v33-activations}
  \centering\scriptsize
  \begin{tabular}{lrrl}
    \toprule
    Mechanism & Units & Activations & Detail \\
    \midrule
    Contract fallback & 0 & 0 & of 1,838 decisions \\
    Bob recovery override & 38 & 46 & 29 offers; 17 accepts \\
    Cycle intervention & 3 & 3 & 1 match; 2 walkaways \\
    Seller short-horizon repair & 3 & 6 & 4 from unknown horizon \\
    Family drawdown stop & 2 & 2 & persuasion; negotiation \\
    Assignment overshoot & 0 & 0 & across 248 checkpoints \\
    Telemetry abort & 0 & 0 & 500 errors after checks \\
    \bottomrule
  \end{tabular}
\end{table}

\paragraph{Telemetry defect.}
The session report contains 500 caught HTTP 422 reward-harvest failures from six
non-UUID component-test IDs. Two IDs were queried 84 times and four were queried
83 times. The ELP wrapper overwrote three synthetic flags, three other IDs did
not match the synthetic prefix classifier, and the checkpoint measured new
telemetry before harvesting. Hence the abort guard saw zero newly appended
errors. The queries occurred after actions and did not change recorded ratings,
but they invalidate any claim that the monitoring layer worked as intended.
The future protocol must preserve synthetic flags, classify test IDs by an
explicit registry, and measure telemetry after harvesting.

\begin{table}[t]
\caption{Descriptive distribution of assignment-level displayed rating changes. No inferential interval is reported because matchmaking, opponent repetition, drift, and stopping are uncontrolled.}
\label{tab:v33-distribution}
\centering\scriptsize
\begin{tabular}{lrrrrr}
\toprule
Family & $+ / 0 / -$ & Mean & Median & SD & IQR \\
\midrule
Bargaining & 76/1/73 & 1.27 & 0.19 & 4.36 & 7.64 \\
Persuasion & 38/0/40 & 0.71 & -0.22 & 3.90 & 6.62 \\
Negotiation & 7/0/13 & -1.69 & -2.98 & 3.65 & 6.16 \\
\bottomrule
\end{tabular}
\end{table}

\begin{table}[t]
\caption{Persuasion role--mode--information support. Hidden-value assignments contain no ELP buyers, so hidden/known marginal differences are not role-adjusted comparisons.}
\label{tab:v33-persuasion-support}
\centering\scriptsize
\begin{tabular}{llllrrr}
\toprule
Role & Mode & Values & $n$ & Sum & Mean & $+ / 0 / -$ \\
\midrule
Buyer & binary & known & 16 & 11.63 & 0.73 & 8/0/8 \\
Buyer & text & known & 21 & 23.14 & 1.10 & 14/0/7 \\
Seller & binary & hidden & 11 & -2.07 & -0.19 & 3/0/8 \\
Seller & binary & known & 9 & 4.79 & 0.53 & 4/0/5 \\
Seller & text & hidden & 15 & -3.93 & -0.26 & 5/0/10 \\
Seller & text & known & 6 & 22.11 & 3.68 & 4/0/2 \\
\bottomrule
\end{tabular}
\end{table}

\begin{table}[t]
\caption{Role-stratified assignment-level rating changes.}
\label{tab:v33-role-distribution}
\centering\scriptsize
\begin{tabular}{llrrrrr}
\toprule
Family & Role & $n$ & Sum & Median & SD & $+ / 0 / -$ \\
\midrule
Bargaining & Alice & 64 & 80.15 & 0.25 & 3.82 & 33/1/30 \\
Bargaining & Bob & 86 & 110.79 & -0.08 & 4.74 & 43/0/43 \\
Persuasion & Buyer & 37 & 34.77 & 0.73 & 3.89 & 22/0/15 \\
Persuasion & Seller & 41 & 20.90 & -0.45 & 3.95 & 16/0/25 \\
Negotiation & Buyer & 11 & -20.33 & -2.91 & 4.02 & 4/0/7 \\
Negotiation & Seller & 9 & -13.40 & -3.04 & 3.37 & 3/0/6 \\
\bottomrule
\end{tabular}
\end{table}

\begin{table}[t]
\caption{Outcome counts in the transcript export. Coverage is 220 of 248 rated assignments, so these are available-case behavioral summaries and are not joined to rating rows.}
\label{tab:v33-transcript-outcomes}
\centering\scriptsize
\begin{tabular}{llrrrr}
\toprule
Family & Role & $n$ & Agreement & Walkaway & No deal \\
\midrule
Bargaining & alice & 54 & 47 & 7 & 0 \\
Bargaining & bob & 73 & 66 & 7 & 0 \\
Persuasion & buyer & 34 & 0 & 0 & 0 \\
Persuasion & seller & 40 & 0 & 0 & 0 \\
Negotiation & buyer & 11 & 6 & 3 & 2 \\
Negotiation & seller & 8 & 5 & 1 & 2 \\
\bottomrule
\end{tabular}
\end{table}

\begin{table}[t]
\caption{Stop-guard sensitivity when all four loss/drawdown thresholds are multiplied by a common factor. Entries are simulated stop game and reason from the observed family sequence.}
\label{tab:v33-stop-sensitivity}
\centering\scriptsize
\begin{tabular}{lrrrrr}
\toprule
Family & $0.50\times$ & $0.75\times$ & $1.00\times$ & $1.25\times$ & $1.50\times$ \\
\midrule
Bargaining & 17 (Bob loss) & 17 (Bob DD) & 150 (target) & 150 (target) & 150 (target) \\
Persuasion & 24 (buyer DD) & 75 (buyer DD) & 78 (family DD) & 78 (end) & 78 (end) \\
Negotiation & 19 (buyer loss) & 20 (buyer DD) & 20 (family DD) & 20 (end) & 20 (end) \\
\bottomrule
\end{tabular}
\end{table}

